\documentclass[10pt,a4paper]{article}
\usepackage[margin=1.8cm]{geometry}
\usepackage[T1]{fontenc}
\usepackage{lmodern}
\usepackage{microtype}
\usepackage{booktabs}
\usepackage{graphicx}
\usepackage{xcolor}
\usepackage{caption}
\usepackage{float}
\usepackage{subcaption}
\usepackage{amsmath}
\usepackage{enumitem}
\usepackage[hidelinks]{hyperref}

\captionsetup{font=small,labelfont=bf}
\setlist{nosep,leftmargin=*}
\setlength{\parskip}{3pt}
\setlength{\parindent}{0pt}

% ---------------------------------------------------------------------------

\newif\ifdraft \draftfalse
\ifdraft
  \newcommand{\chk}[1]{\textcolor{orange!85!black}{#1}}
  \newcommand{\tbd}{\textcolor{red}{\textbf{??}}}
\else
  \newcommand{\chk}[1]{#1}
  \newcommand{\tbd}{--}
\fi
% Figure that shows a placeholder box until the PNG exists in figures/
\newcommand{\fig}[2]{%
  \IfFileExists{figures/#1}{\includegraphics[width=#2]{figures/#1}}%
  {\fbox{\parbox[c][0.32\linewidth][c]{#2}{\centering\small\texttt{\detokenize{figures/#1}}\par(pending)}}}}
% ---------------------------------------------------------------------------

\begin{document}

{\Large\bfseries GEN4 at Jarama: lap time and energy projection\\[2pt]
from a Gen3 Evo lap simulation calibrated on the 2026 Madrid E-Prix}

\smallskip
Carlos S\'anchez Maceda \hfill 9 November 2026 \hfill
\href{https://github.com/cjsanchezmaceda/FE_Energy_LapSim}{\texttt{github.com/cjsanchezmaceda/FE\_Energy\_LapSim}}

\noindent\rule{\linewidth}{0.4pt}

\textbf{Summary.}
A quasi-steady-state (QSS) lap simulation of the Formula~E layout at Jarama was calibrated on the
best dry FP1 lap of the 2026 Madrid E-Prix and used to project the GEN4 car ahead of the
16--20 November pre-season test.
\begin{itemize}
  \item \textbf{Gen3 Evo calibration:} one free parameter (tyre $\mu = 1.689$) reproduces the
        350\,kW lap (89.403\,s); the independent 300\,kW lap is predicted to within
        +0.14\,s. The result is insensitive to the unknown drag area and mass.
  \item \textbf{GEN4, unconstrained:} \textbf{80.6\,s} in qualifying trim (600\,kW, high-DF),
        $-8.8$\,s vs Gen3 Evo; high-DF is 1.7\,s faster than low-DF at 600\,kW.
  \item \textbf{GEN4, energy-limited (450\,kW):} the model supports the intended use of the two kits
        (high-DF for qualifying, low-DF for races): \textbf{low-DF is faster for any race budget below
        $\approx$3.6\,kWh/lap}, i.e.\ any race longer than 14 laps on the 51.25\,kWh usable energy;
        by 3.5\,s at 2.6\,kWh/lap. Lift-and-coast meets the same budget
        0.9--3.6\,s/lap faster than capping drive power (2.4\,s for Gen3 Evo).
\end{itemize}

% ===========================================================================
\section*{1\quad Model}

\textbf{Track.} Track limits digitised from satellite imagery (Google Earth, $\pm0.6$\,m vs GPS
boundaries), including the FE loop (+30.1\,m). Centreline 3933\,m vs 3934\,m official.
Racing line: minimum-curvature path (Heilmeier et al., 2020) with 1.3\,m margin to the limits,
3882\,m, minimum radius 11.6\,m; curvature from a periodic cubic spline, 2\,m resolution.

\textbf{Lap simulation.} Point-mass QSS: cornering limit
$v^2 = \mu g / (|\kappa| - \mu\rho C_LA/2m)$, backward (braking) and forward (traction) passes on a
two-lap domain. Combined grip via the friction circle; longitudinal load transfer
($F h/L$) in traction and braking; battery-side power limits with drive/regen efficiencies;
Gen3 rear axle braked by regen only (350\,kW), GEN4 with rear friction brakes.
Energy: $E = \sum P\,\Delta t$ from the wheel power, net of regeneration.

\begin{table}[H]
\centering\small
\caption{Vehicle parameters. \emph{Public}: FIA regulations / Formula~E. \emph{Assumed}: not public, quoted with the range tested.}
\begin{tabular}{@{}llll@{}}
\toprule
 & Gen3 Evo & GEN4 & Source \\
\midrule
Mass incl.\ driver [kg]        & 862                & 1046                 & FIA Tech.\ Regs (public) \\
Wheelbase / width [m]          & 2.97 / 1.70        & 3.08 / 1.79          & FIA Tech.\ Regs (public) \\
Usable race energy [kWh]       & 38.5 + 3.85 PIT BOOST & 51.25             & public \\
Drive power [kW]               & 300 race / 350     & 450 race / 600       & public \\
Driven axles                   & RWD; AWD at 350\,kW & AWD                 & public \\
Regen [kW] front / rear        & 250 / 350          & 350 / 350            & FIA Tech.\ Regs (public) \\
Rear friction brakes           & no                 & yes                  & public \\
Tyre $\mu$                     & 1.689 (calibrated) & $1.689\times1.05$     & tyre factor assumed, 1.00--1.10 \\
$C_DA$ [m$^2$] low / high-DF   & 0.90               & 0.95 / 1.24          & assumed; high-DF $\approx$+30\,\% drag \\
$C_LA$ [m$^2$] low / high-DF   & 1.6                & 2.7 / 3.8            & assumed \\
Front motor [kW]               & 50                 & 250                  & FIA Tech.\ Regs (public) \\
$\eta$ drive / regen           & 0.93 / 0.90        & 0.93 / 0.90          & assumed \\
\bottomrule
\end{tabular}
\end{table}

% ===========================================================================
\section*{2\quad Gen3 Evo calibration (Madrid E-Prix 2026, FP1, dry)}

Only the tyre friction coefficient is fitted. Target: best 350\,kW lap of FP1 (Ticktum,
1:29.403, best of 16 drivers who ran 350\,kW). Check: best 300\,kW lap (Da Costa, 1:31.612,
best of only 4 drivers, hence an upper bound rather than a limit lap).

\begin{table}[H]
\centering\small
\caption{Calibration and robustness. Each row re-fits $\mu$ on the 350\,kW lap.}
\begin{tabular}{@{}lccc@{}}
\toprule
Case & $\mu$ & 300\,kW lap [s] & vs FIA 91.612\,s \\
\midrule
\textbf{Nominal} ($C_DA$ 0.90, AWD at 350\,kW, 862\,kg) & \textbf{1.689} & \textbf{91.76} & \textbf{+0.14} \\
$C_DA$ 0.75 / 1.05                                     & 1.675 / 1.703 & 91.70 / 91.83 & +0.08 / +0.22 \\
Mass 840 / 936\,kg                                     & 1.674 / 1.737 & 91.77 / 91.73 & +0.16 / +0.12 \\
350\,kW in FP1 without front motor (RWD)               & 1.734         & 90.83         & $-0.78$ \\
\bottomrule
\end{tabular}
\end{table}

$C_DA$ and mass are absorbed by $\mu$ with $\le 0.07$\,s effect on the prediction. The main open question is whether the front motor was active on the 350\,kW FP1 laps.
If it was (nominal), the model reproduces the FIA gap between the 350 and 300\,kW laps
(2.36 vs 2.21\,s). If it was not, $\mu$ rises to 1.734 and the model would place the
best 300\,kW lap 0.8\,s slower than the car's limit, which is less plausible.

% ===========================================================================
\section*{3\quad GEN4 projection, unconstrained lap}

\begin{table}[H]
\centering\small
\caption{GEN4 flat-out laps, FE layout. $\Delta$ vs Gen3 Evo 350\,kW (89.40\,s).}
\begin{tabular}{@{}lcccccc@{}}
\toprule
Kit & Power [kW] & Lap [s] & $\Delta$ [s] & $v_{max}$ [km/h] & $v_{min}$ [km/h] & Energy [kWh] \\
\midrule
low-DF  & 450 & 84.97 & $-4.42$ & 273 & 52 & 3.61 \\
low-DF  & 600 & 82.26 & $-7.14$ & 293 & 52 & 4.36 \\
high-DF & 450 & 83.65 & $-5.74$ & 267 & 52 & 4.37 \\
\textbf{high-DF} & \textbf{600} & \textbf{80.59} & $\mathbf{-8.80}$ & 290 & 52 & 5.21 \\
\bottomrule
\end{tabular}
\end{table}

\begin{figure}[H]
\centering
\fig{fig1_speed_gen3_gen4.png}{0.95\linewidth}
\caption{Speed vs distance: Gen3 Evo 350\,kW (calibrated) vs GEN4 600\,kW high-DF. Minimum speeds
in the slow corners are almost unchanged (52 vs 50\,km/h): the GEN4 gains come from power on
the straights (+25--40\,km/h at the end of each, top speed 290 vs 256\,km/h) and from downforce in the medium-speed corners.}
\end{figure}

% ===========================================================================
\section*{4\quad Energy-limited laps}

Two ways of meeting an energy budget per lap. \emph{Lift-and-coast}: the driver lifts off over
the last fraction $f_c$ of each acceleration phase, measured in distance back from the flat-out
braking point (no drive, no regen). \emph{Power cap}: lower maximum drive power over the whole lap.
$f_c$ or the cap is solved so the lap uses exactly the budget.

\textbf{Gen3 Evo check vs the race.} Budget $(38.5+3.85)/23 = 1.84$\,kWh/lap (usable energy +
PIT BOOST, 23 laps).

\begin{table}[H]
\centering\small
\caption{Gen3 Evo, 300\,kW, Madrid race energy budget.}
\begin{tabular}{@{}lccc@{}}
\toprule
Strategy & Lap [s] & Energy [kWh] & Setting \\
\midrule
Flat-out (no budget)   & 91.75 & 2.34 & -- \\
Lift-and-coast         & \textbf{95.18} & 1.84 & $f_c = 43$\,\% \\
Power cap              & 97.55 & 1.84 & 186\,kW \\
Race: best normal laps, 4 of 5 drivers & 94.8--95.7 & -- & excl.\ lap 1, pit, ATTACK MODE \\
Race: median normal lap & 98.2 & -- & 80 laps, 5 drivers \\
\bottomrule
\end{tabular}
\end{table}

The energy-limited lift-and-coast lap (95.18\,s) falls within the best normal race laps of four
of the five drivers transcribed, and 3.0\,s below the median, consistent with traffic and racing
effects that the model does not include. Power capping would cost 2.4\,s/lap more.

\begin{figure}[H]
\centering
\fig{fig2_energy_tradeoff_gen4.png}{0.6\linewidth}
\caption{GEN4 at 450\,kW: lap time vs battery energy per lap, by aero kit and saving strategy.
The kits cross at 3.60\,kWh/lap: below it, low-DF is faster. On the 51.25\,kWh GEN4 usable
energy, 3.6\,kWh/lap corresponds to a 14-lap race.}
\end{figure}

\begin{table}[H]
\centering\small
\caption{GEN4 at 450\,kW, lap time [s] for a given energy per lap. Low-DF needs no saving above its flat-out 3.61\,kWh.}
\begin{tabular}{@{}lccccc@{}}
\toprule
Energy [kWh/lap] & 2.6 & 3.0 & 3.4 & 3.8 & flat-out \\
\midrule
low-DF        & 88.66 & 86.15 & 85.08 & 84.97 & 84.97 \\
high-DF       & 92.14 & 88.34 & 85.81 & 84.39 & 83.65 \\
$\Delta$ (high $-$ low) & $+3.48$ & $+2.19$ & $+0.73$ & $-0.58$ & $-1.32$ \\
\midrule
Power cap, low-DF  & 91.90 & 88.48 & 85.98 & 84.97 & 84.97 \\
Power cap, high-DF & 95.53 & 91.90 & 88.86 & 86.42 & 83.65 \\
\bottomrule
\end{tabular}
\end{table}

% ===========================================================================
\section*{5\quad Does the kit conclusion survive the assumptions?}

\begin{table}[H]
\centering\small
\caption{$\Delta$ lap time high-DF $-$ low-DF [s], one assumption changed at a time.
Negative: high-DF faster.}
\begin{tabular}{@{}lcc@{}}
\toprule
Case & flat-out, 450\,kW & 3.0\,kWh/lap, lift-and-coast \\
\midrule
Nominal                         & $-1.32$ & $+2.20$ \\
Tyre factor 1.00 / 1.10         & $-1.38$ / $-1.27$ & $+1.80$ / $+2.09$ \\
High-DF $C_LA$ $-$15 / +15\,\%  & $-0.36$ / $-2.21$ & $+2.34$ / $+1.77$ \\
High-DF $C_DA$ $-$10 / +10\,\%  & $-1.59$ / $-1.06$ & $+0.92$ / $+3.52$ \\
\bottomrule
\end{tabular}
\end{table}

The sign of the conclusion holds in every case: high-DF is faster on an unconstrained lap,
low-DF is faster once energy is limited. The high-DF drag area is the assumption that matters
most for the race case (0.9--3.5\,s); downforce matters most for the flat-out lap.
The GEN4 test speed traps would pin $C_DA$ down directly.

% ===========================================================================
\section*{6\quad Limitations and next steps}

\textbf{Limitations.} Point-mass QSS (no transients, no tyre temperature); single $\mu$ for both
axles; front-motor traction uses static front load; lift-and-coast is a fixed fraction per
acceleration phase, not an optimal energy allocation; drivetrain efficiencies and all GEN4 aero
and tyre values are assumed; GEN4 regen split assumed; whether 350\,kW FP1 running used the
front motor is unconfirmed; track digitised, not surveyed.

\textbf{Next steps.}
(1)~Re-fit tyre factor and $C_DA$ from GEN4 test laps and speed traps; the model then becomes a
GEN4 tool.
(2)~Optimal-control energy allocation (where to coast, lap by lap) instead of a uniform fraction.
(3)~ATTACK MODE and PIT BOOST deployment strategy over a race distance.
(4)~Slipstream: energy saved per lap when following.

\end{document}